\hmsubsection{System Reliability}{FKT}{}

System reliability was assessed through subsystem testing followed
by the final integrated MVP validation session. Earlier subsystem
work verified the camera pipeline, inverse kinematics solver,
OpenRB-150 communication, motor actuation, calibration routines,
and adaptive grip verification. The final validation session then
tested whether these subsystems operated together during repeated
scan-pick-place cycles with real hardware.

The console log for the final session confirmed that the system
used a real serial connection to the OpenRB-150 and a real OAK-D
camera. It also confirmed that the inverse kinematics solver was
ready, the OpenRB handshake was completed, and adaptive grip
verification was active before repeated scan-pick-place execution.
This is important because the reported reliability result is based
on the integrated physical system rather than simulation or mock
interfaces.

In the final 20-ball validation session, the system sorted
\textbf{20/20 balls correctly (100\% observed net sorting
success)} and recorded \textbf{0 unrecovered failures}. The system
picked and placed 10/10 blue balls and 10/10 red balls correctly.
One additional grip-verification rejection occurred in cycle~7 on
a red ball before successful retry. Grip verification correctly
rejected an unstable grasp, and retry logic recovered the cycle on
the next scan. This event therefore demonstrates the usefulness of
the recovery mechanism, but it is reported separately from the net
sorting result and is not counted as an unrecovered sorting
failure.

The reliability conclusion is consequently limited and specific:
the MVP was achieved under the tested laboratory conditions in the
final validation session. The result does not establish perfect
general reliability for all lighting conditions, object positions,
or long-duration operation. It does, however, provide evidence
that the integrated system can complete the required autonomous
sorting task with no unrecovered failures in the tested 20-cycle
session.
